\ncdlesson{1}{Hvad mener danskerne egentlig, når de siger "det er fint"?}{丹麦人说“挺好”时到底是什么意思？}
{DU3 Modul 4（B1+ $\rightarrow$ B2-）｜对标 NCE 3 \& Modultest 4 \textit{Læsning Opgave 1--2 \& Skrivning Opgave 1 (Email)}}
{\textbf{灵魂情态小品词 I}：\textit{jo}（你我皆知/共识）· \textit{da}（当然/略带反驳/惊叹）· \textit{nok}（大概/恐怕/安心预测）· \textit{vist}（听说/好像）· 情态小品词在主句与从句中的\textbf{中场语序定位}}
{欢迎来到《新概念丹麦语 4 · 思辨篇》！从 B1 迈向 B1+/B2-，决定地道程度的分水岭往往不在高深长词，而在看似不起眼的“情态小品词（\textit{modalpartikler}）”。丹麦人说话从不直接把话说死，常把真正立场藏在 \textit{jo, da, nok, vist} 里：少一个词语气生硬，用错一个词意味全反。吃透本课，你不仅能听懂丹麦同事的“弦外之音”，写 Modul 4 私人长邮件时也能瞬间摆脱僵硬的教科书腔！}

% ---------------------------------------------------------------------
\begin{foerduleaeser}
带着下面两个问题阅读 Tekst A 和 Tekst B：
\begin{enumerate}
  \item \textit{Hvorfor misforstår Emma først Henriks kommentar ”det er vist meget fint” i Tekst A, og hvordan forklarer Sofie betydningen af ”vist” og ”jo”?}（在 Tekst A 中，艾玛最初为什么误解了亨里克那句“det er vist meget fint”？索菲是如何解释 \textit{vist} 和 \textit{jo} 的含义的？）
  \item \textit{Hvorfor er Kasper bekymret før fredagsmødet i Tekst B, og hvordan bruger Freja ”jo”, ”da” og ”nok” til at berolige ham?}（在 Tekst B 中，卡斯珀为什么在周五会议前感到担忧？弗蕾娅如何用 \textit{jo}、\textit{da} 和 \textit{nok} 来宽慰他？）
\end{enumerate}
\end{foerduleaeser}

% ---------------------------------------------------------------------
\ncdsubsec{二、Tekster（口语微调 / 报刊与思辨双语篇）}

\subsection*{Tekst A：\textit{Den skjulte betydning af ”det er fint”}（职场文化随笔 · 约 235 词）}

\begin{tipbox}{语篇说明}
本篇聚焦丹麦职场中极易引发跨文化误解的“含蓄表达”。阅读时请紧盯两种色彩标记：\textbf{\color{NCDNavy}深蓝色加粗的 \textit{jo / da}}（诉诸共识或温和反驳/鼓励）与\textbf{\color{NCDRed}红色加粗的 \textit{nok / vist}}（推测、宽慰或委婉存疑）。
\end{tipbox}

Da Emma flyttede til Aarhus og begyndte i en dansk virksomhed, troede hun, at hun hurtigt forstod sine kolleger. Efter sin første store præsentation spurgte hun sin kollega Henrik, hvad han syntes om hendes udkast. Henrik smilte roligt og svarede: ”Det er \dabold{vist} meget fint, men vi kan \da{jo} lige kigge på tidsplanen i morgen.”

Emma blev lettet og troede, at Henrik var helt overbevist. Ved kaffemaskinen forklarede hendes veninde Sofie dog, at danske småord ofte gemmer på vigtige nuancer. ”Når Henrik siger \dabold{vist}, betyder det faktisk, at han ikke er helt sikker, eller at han tvivler en smule,” sagde Sofie. ”Og med \da{jo} minder han dig høfligt om noget, som I begge to godt ved --- nemlig at tidsplanen stadig mangler arbejde.”

Efterhånden opdagede Emma, hvor stor forskel disse fire småord gør i hverdagen. Bruger en dansker \da{da}, ligger der ofte en mild indvending, forbavselse eller opmuntring i tonen, som i: ”Du kommer \da{da} til mødet, ikke?” eller ”Det var \da{da} ærgerligt!” Siger kollegaen derimod \dabold{nok}, udtrykker det enten en beroligende formodning eller en forsigtig indrømmelse: ”Det går \dabold{nok} godt,” eller ”Du har \dabold{nok} ret, men vi bør vente lidt.”

Samtidig må man passe på ordstillingen. I hovedsætninger står småordene efter det første verbum, mens de i ledsætninger sniger sig ind foran verbet, ligesom \textit{ikke}. Når man først lærer disse nuancer at kende, bliver danskernes sprog pludselig langt lettere at afkode.

\begin{transbox}
\small
\textbf{\color{NCDTeal}[Tekst A 中文精译]}\\
当艾玛搬到奥胡斯并在一家丹麦企业开始工作时，她原本以为自己很快就能听懂同事们的话。在完成第一次大型汇报后，她问同事亨里克对自己的方案初稿怎么看。亨里克平静地微笑着回答：“这大概挺不错的吧，不过我们明儿当然可以顺便再看一下时间表。”

艾玛松了一口气，以为亨里克已经完全被说服了。然而在咖啡机旁，她的好友索菲却解释说，丹麦语里的小品词常常藏着关键的语气微差。“当亨里克说 \textit{vist} 的时候，其实意味着他还不太确定，或者稍微有点疑虑，”索菲说道，“而用 \textit{jo}，他是在礼貌地提醒你一件你们俩心知肚明的事——那就是时间表还需要再打磨。”

渐渐地，艾玛发现这四个小词在日常生活中能带来天壤之别。当丹麦人使用 \textit{da} 时，语气里往往带有一丝温和的反驳、惊讶或鼓励，比如：“你肯定会来开会的吧？”或者“那可真是太遗憾了！”相反，如果同事说 \textit{nok}，通常表达一种令人安心的推测或谨慎的让步：“一切大概会顺利的”，或者“你说得恐怕没错，但我们还是应该再等等”。

与此同时，人还必须留心语序。在主句中，这些小词位于第一个变位动词之后；而在从句中，它们会像 \textit{ikke} 一样溜到变位动词前面。一旦掌握了这些细微差别，丹麦人说话背后的真实意图一下子就变得清晰可读了。
\end{transbox}

\subsection*{Tekst B：\textit{Intern mail: Nervøs før fredagsmødet}（对标 DU3 Modul 4 私人与职场回信 · 约 130 词）}

\begin{emailbox}{来信｜Fra: Kasper \quad Til: Freja \quad Emne: Lidt nervøs før præsentationen}
\textbf{Hej Freja}\\
Jeg skal præsentere vores nye budget på fredag, og jeg er ærlig talt lidt bekymret. Sidst jeg sendte et forslag til direktøren, skrev han bare, at det \dabold{vist} så fint ud, og så hørte jeg ikke mere. Tror du, at han var skeptisk?\\
\textbf{Venlig hilsen, Kasper}
\end{emailbox}

\begin{emailbox}{回复｜Fra: Freja \quad Til: Kasper \quad Emne: Sv: Lidt nervøs før præsentationen}
\textbf{Hej Kasper}\\
Tag det helt roligt! Det betyder \da{jo} ikke automatisk, at han var utilfreds --- han havde \dabold{vist} bare utrolig travlt den uge. Du ved \da{jo}, at vores chef sætter pris på konkrete tal, og dine analyser er \da{da} altid grundige og klare! Hvis vi øver os en halv time i morgen, skal det \dabold{nok} gå helt fint på fredag. Vi tales ved!\\
\textbf{Bedste hilsner, Freja}
\end{emailbox}

\textbf{Tekst B 阅读核对（Rigtigt / Forkert / Kort svar，答案见第六部分）：}
\begin{enumerate}
  \item \textit{Direktøren gav Kasper en lang og meget begejstret tilbagemelding sidste gang.} \underline{\hspace{2.2cm}}
  \item \textit{Hvordan forklarer Freja, at direktøren ikke skrev mere efter det sidste forslag?} \\\underline{\hspace{5.2cm}}
  \item \textit{Hvad foreslår Freja at gøre i morgen, og hvilken sikkerhed udtrykker hun med ”skal det nok gå”?} \underline{\hspace{5.2cm}}
\end{enumerate}

\begin{tipbox}{本课发音与语调要点}
\begin{itemize}
  \item \textbf{轻读 vs 重读决定词性}：情态小品词在句中一律\textbf{轻读且不带喉塞音（\textit{stød}）}：\textit{jo} \ipa{jo}、\textit{da} \ipa{da}、\textit{nok} \ipa{nʌɡ}、\textit{vist} \ipa{vesd}。一旦重读，词义立刻改变：\textit{\textbf{Jo!}} \ipa{ˈjoːˀ}（用于反驳否定问句的“不，恰恰相反/是的”）、\textit{\textbf{da}} \ipa{ˈdaː}（“那时”）、\textit{\textbf{nok}} \ipa{ˈnʌɡ}（“足够”，如 \textit{tid nok}）。
  \item \textbf{极易忽略的静音字母}：\textit{egentlig} \ipa{ˈeːˀənli}（\textbf{g 和 t 均不发音}）；\textit{indvending} \ipa{ˈenˌvɛneŋ} 与 \textit{indrømmelse} \ipa{ˈenˌʁømˀəlsə}（前缀 \textit{ind-} 里的 \textbf{d 不发音}）；\textit{ærgerligt} \ipa{ˈɛɐ̯ˀɐli}（\textbf{g 和 t 均不发音}！）。
  \item \textbf{核心词重音与软 d}：\textit{formodning} \ipa{fʌˈmoːˀðneŋ}（重音在第二音节，含软 d）；\textit{overbevist} \ipa{ˈʌʊ̯ɐbeˌviːˀsd}；\textit{tidsplan} \ipa{ˈtˢiðsˌpʰlæːˀn}。
\end{itemize}
\end{tipbox}

\begin{transbox}
\small
\textbf{\color{NCDTeal}[Tekst B 中文对照]}\\
\textbf{卡斯珀的邮件}：你好，弗蕾娅。我周五要汇报我们的新预算，老实说有点担心。上次我把一份提案发给总经理时，他只回了一句说看上去大概还行，之后就没下文了。你觉得他当时是持怀疑态度吗？诚挚的问候，卡斯珀。

\textbf{弗蕾娅的回复}：你好，卡斯珀。完全放轻松！这当然并不自动意味着他不满意——他那周大概只是忙得不可开交罢了。你本来也知道我们老板最看重扎实的数据，而且你的分析本来向来就非常细致清晰嘛！如果我们明天抽出半小时一起排练一下，周五肯定没问题的。回头聊！谨致问候，弗蕾娅。
\end{transbox}

% ---------------------------------------------------------------------
\ncdsubsec{三、Gloser（核心生词表）}

\begin{tipbox}{词表重点}
\textit{jo, da, nok, vist} 作为情态小品词时均无词形变化，但其语义高度依赖上下文语境与重音位置。名词 \textit{en nuance} 与 \textit{en indrømmelse}（让步）是 Modul 4 讨论语气、写作审题时的必备核心词；动词 \textit{at tvivle} 固定搭配介词 \textit{på}（怀疑某事）。
\end{tipbox}

\begin{glosertable}
小品词 & jo & \textit{(不变，情态小品词)} & 你我皆知的，毕竟，本来嘛 \\
小品词 & da & \textit{(不变，情态小品词)} & 当然，难道不是吗，可真是 \\
小品词 & nok & \textit{(不变，情态小品词)} & 大概，恐怕，放心会……的 \\
小品词 & vist & \textit{(不变，情态小品词)} & 听说，好像，大概（含不确定） \\
\midrule
名词 & en nuance & \textit{(nuancen, nuancer, nuancerne)} & 细微差别，语气微差 \\
名词 & et udkast & \textit{(udkastet, udkast, udkastene)} & 草稿，初稿（单复同形） \\
名词 & en tidsplan & \textit{(tidsplanen, tidsplaner, -erne)} & 时间表，进度计划 \\
名词 & en indvending & \textit{(indvendingen, indvendinger, -erne)} & 异议，反对意见 \\
名词 & en formodning & \textit{(formodningen, formodninger, -erne)} & 推测，假设，猜想 \\
名词 & en indrømmelse & \textit{(indrømmelsen, indrømmelser, -erne)} & 让步，承认 \\
名词 & en bekymring & \textit{(bekymringen, bekymringer, -erne)} & 担忧，顾虑 \\
\midrule
动词 [har] & at mene & \textit{(mener, mente, har ment)} & 认为，意思是指，意欲表达 \\
动词 [har] & at tvivle (på) & \textit{(tvivler, tvivlede, har tvivlet)} & 怀疑，拿不准 \\
动词 [har] & at berolige & \textit{(beroliger, beroligede, har beroliget)} & 安抚，使安心 \\
动词 [har] & at afkode & \textit{(afkoder, afkodede, har afkodet)} & 解码，领会弦外之音 \\
\midrule
形容词 & overbevist & \textit{(overbevist, overbeviste)} & 信服的，确信的 \\
形容词 & skeptisk & \textit{(skeptisk, skeptiske)} & 持怀疑态度的 \\
\end{glosertable}

% ---------------------------------------------------------------------
\ncdsubsec{四、Grammatik \& Stilistik（核心语法、语气与修辞精讲）}

\begin{grammatikbox}{语法点 1：四大灵魂情态小品词（\textit{jo, da, nok, vist}）语用潜台词全景对照}
这四个单音节词在英文里没有直接一一对应的词，因此常被初学者忽略，却决定着句子听起来像“冷冰冰的机器人”还是“懂分寸的丹麦人”：

\vspace{0.3em}
\renewcommand{\arraystretch}{1.25}
{\small
\begin{tabularx}{\linewidth}{@{} >{\bfseries\raggedright\arraybackslash}p{1.6cm} >{\raggedright\arraybackslash}p{3.6cm} >{\raggedright\arraybackslash}p{4.3cm} >{\raggedright\arraybackslash}X @{}}
\toprule
\textbf{小品词} & \textbf{核心语用功能} & \textbf{潜台词与中文神韵} & \textbf{典型例句与语气解析} \\
\midrule
\da{jo} & 诉诸双方共识\newline （提醒已知事实） & “你我都知道的呀 / 毕竟……嘛”（暗含：这事没必要再争了） & \textit{Vi har \da{jo} aftalt mødet i går.}（我们昨天毕竟约好开会了呀。）\newline \textit{Det tager \da{jo} tid at lære dansk.} \\
\midrule
\da{da} & 温和反驳、催促、安慰或惊叹 & “当然啦 / 可不是嘛 / 怎么不……呢？”（带情绪温度与互动感） & \textit{Det var \da{da} en god idé!}（这主意可真不错！）\newline \textit{Du er \da{da} ikke sur?}（你该不会生气了吧？） \\
\midrule
\dabold{nok} & 说话人自身较有把握的推测、宽慰或让步 & “大概会…… / 放心吧会…… / 恐怕确实如此，但……” & \textit{Han skal \dabold{nok} nå toget.}（他肯定赶得上火车，放心。）\newline \textit{Det er \dabold{nok} dyrt, men kvaliteten er høj.} \\
\midrule
\dabold{vist} & 转述听说、模糊记忆或委婉存疑 & “听说…… / 我好像记得…… / 恐怕有点……（留余地）” & \textit{Han kommer \dabold{vist} først i morgen.}（他好像明天才来，我不太确定。）\newline \textit{Mødet var \dabold{vist} lidt langt.} \\
\bottomrule
\end{tabularx}
}
\vspace{0.3em}

\textbf{高频对比辨析}：\ding{172} \dabold{nok} vs \dabold{vist}：说 \textit{Det går \dabold{nok} godt} 是安慰对方“肯定没问题”；说 \textit{Det gik \dabold{vist} meget godt} 则是“听说还行吧（但说话人自己心里可能没底或有所保留）”。\ding{173} \da{jo} vs \da{da}：\textit{Du ved \da{jo}, at...} 强调“你本来就知道”；\textit{Du ved \da{da}, at...} 往往在反问或反驳对方忘了这件事（“你明明知道的呀！”）。
\end{grammatikbox}

\begin{grammatikbox}{语法点 2：语序铁律——情态小品词在句子“中场（Midterfeltet）”的精确落点}
情态小品词与否定词 \textit{ikke}、时间副词 \textit{altid / ofte} 同属于丹麦语拓扑句法中的\textbf{中场副词（\textit{sætningsadverbialer}）}，位置错一步整句就变味：
\begin{enumerate}
  \item \textbf{主句（Hovedsætning）中的 V2 定位}：\ding{172} 若主语在句首，小品词紧跟在\textbf{第一个变位动词}之后：\textit{Hun kommer \da{nok} i dag.}（不可以说 \ding{55}\ \textit{Hun nok kommer}）。\ding{173} 若时间/地点或从句占领句首触发 V2 倒装，小品词放在\textbf{动词 + 主语}之后：\textit{I morgen kommer hun \da{nok}.}
  \item \textbf{遇上轻代词宾语与 \textit{ikke} 谁先谁后？}：\ding{174} 没有重读的短代词宾语（\textit{det, ham, hende, mig, sig}）必须穿过中场跑到小品词\textbf{前面}：\textit{Jeg sagde \dabold{det jo} til dig!} / \textit{Han husker \dabold{det vist} ikke.} \ding{175} 情态小品词通常排在否定词 \dabold{ikke} \textbf{前面}：\textit{Det er \da{jo ikke} så svært.} / \textit{Hun kommer \dabold{nok ikke} til festen.}
  \item \textbf{从句（Ledsætning）语序铁律}：在 \textit{at, fordi, når, hvis, som} 等引导的从句中，小品词必须紧贴主语之后、跑到\textbf{变位动词之前}：
    \begin{itemize}
      \item \ding{51}\ \textit{Han blev hjemme, fordi han \dabold{jo} havde travlt.}（\ding{55}\ \textit{...fordi han havde jo travlt.}）
      \item \ding{51}\ \textit{Hun fortalte, at direktøren \dabold{vist ikke} var helt tilfreds.}
    \end{itemize}
\end{enumerate}
\end{grammatikbox}

\begin{grammatikbox}{地道交际表达与文化贴士（Pragmatik på arbejdspladsen）}
\begin{itemize}
  \item \da{Det skal nok gå!} —— 这事肯定没问题的！（丹麦人安慰朋友、同事的万能金句）。
  \item \da{Det var da ærgerligt / dejligt!} —— 那可真遗憾 / 太棒了！（加 \textit{da} 让情绪共鸣自然真实）。
  \item \textbf{文化提醒}：在丹麦职场听评价时，如果对方用 \textit{vist} 或 \textit{nok} 搭配平淡形容词（如 \textit{Det ser vist okay ud}），多半意味着还有修改空间，千万别误以为对方已经“百分百满意”！
\end{itemize}
\end{grammatikbox}

% ---------------------------------------------------------------------
\ncdsubsec{五、Mønsterøvelser（句型与修辞操练）}

\textbf{练习 A：语序定位手术——将括号中的情态小品词插入句中唯一准确位置（注意 V2 倒装、轻代词与从句语序）}
\begin{enumerate}
  \item \textit{Jeg sagde det til dig i går.}（插入 \textit{jo}）$\rightarrow$
  \item \textit{I morgen når hun ikke at sende rapporten.}（插入 \textit{nok}）$\rightarrow$
  \item \textit{Hun blev overrasket, fordi hun ikke kendte reglerne.}（在从句中插入 \textit{jo}）$\rightarrow$
  \item \textit{De var meget glade for forslaget.}（插入 \textit{da} 表温和感叹/反驳）$\rightarrow$
  \item \textit{Direktøren rejser til Berlin på tirsdag.}（插入 \textit{vist} 表听说/不确定）$\rightarrow$
  \item \textit{Han tror, at mødet begynder klokken ti.}（在从句中插入 \textit{vist}）$\rightarrow$
  \item \textit{Du har glemt din paraply igen.}（插入 \textit{da} 温和提醒）$\rightarrow$
\end{enumerate}

\textbf{练习 B：语用潜台词选填（从 \textit{jo, da, nok, vist} 中选择最符合语境的词填空）}
\begin{enumerate}
  \item Du behøver ikke at være nervøs for eksamen --- det skal \underline{\hspace{2.2cm}} gå rigtig godt!
  \item Hvorfor spørger du mig om adressen? Du har \underline{\hspace{2.2cm}} været hos mig tre gange før!
  \item Jeg kan desværre ikke huske det præcise tal, men det kostede \underline{\hspace{2.2cm}} omkring 500 kroner.
  \item Nå, har du allerede mistet din cykel? Det var \underline{\hspace{2.2cm}} utrolig ærgerligt!
  \item Han lovede at ringe klokken to, men han har \underline{\hspace{2.2cm}} glemt det, fordi han sidder i møde.
  \item Kommer du ikke med til frokost? --- \underline{\hspace{2.2cm}}, selvfølgelig gør jeg det!（注意回答否定问句时的重读词！）
\end{enumerate}

\textbf{练习 C：中译丹实战（活用情态小品词传递潜台词与准确语序）}
\begin{enumerate}
  \item 放心吧，明天他大概不会迟到。（用 \textit{nok}）
  \item 我毕竟已经把邮件发给你了呀。（用 \textit{jo}）
  \item 听说我们的老板下周要去奥胡斯。（用 \textit{vist}）
  \item 你当然可以问同事这个问题！（用 \textit{da}）
  \item 她没来上班，因为她好像生病了。（在从句中用 \textit{vist}）
  \item 他大概是对的，但我们还需要再核对一下数据。（用 \textit{nok}）
\end{enumerate}

% ---------------------------------------------------------------------
\ncdsubsec{六、DU3 Modul 4 实战微任务与参考答案}

\begin{emailbox}{DU3 Modul 4 写作任务（对标 Skrivning Opgave 1 · 私人回复邮件 $\ge$ 125 ord）}
\textbf{Situation}: Din ven Mikkel er lige startet i et nyt job i København. Han har skrevet til dig, fordi hans danske chef sagde: ”Det er vist meget fint”, da Mikkel viste sit første projekt. Nu er Mikkel bange for, at chefen slet ikke kan lide hans arbejde.\\
\textbf{Opgave}: Skriv en beroligende og hjælpsom mail til Mikkel på \textbf{mindst 125 ord} (gerne 140--160 ord). Forklar ham, hvad danske småord som \textit{vist}, \textit{jo}, \textit{nok} og \textit{da} kan betyde i sådan en situation, giv ham et konkret råd til næste møde med chefen, og brug selv alle fire modalpartikler naturligt i din mail!
\end{emailbox}

\begin{answerbox}
\textbf{Før du læser 参考回答}
\begin{enumerate}
  \item \textit{Emma tror først, at Henrik er helt overbevist om hendes udkast, men Sofie forklarer, at ”vist” viser en lille tvivl eller usikkerhed, mens ”jo” høfligt minder om noget, de begge to ved --- at tidsplanen stadig mangler arbejde.}
  \item \textit{Kasper er bekymret, fordi direktøren sidst kun skrev, at hans forslag ”vist så fint ud”. Freja bruger ”jo” til at minde om direktørens travlhed og stil, ”da” til at rose Kaspers analyser og ”nok” til at forsikre ham om, at det går godt.}
\end{enumerate}

\textbf{Tekst B 阅读核对答案}
\begin{enumerate}
  \item \da{Forkert}（总经理只回复了一句 \textit{at det vist så fint ud}，之后没有再多说）
  \item \da{Hun mener, at han vist bare havde utrolig travlt den uge.}
  \item \da{Hun foreslår, at de øver sig en halv time i morgen, og viser tillid til, at det går godt.}
\end{enumerate}

\textbf{练习 A 答案}
\begin{enumerate}
  \item \textit{Jeg sagde det \da{jo} til dig i går.}（轻代词 \textit{det} 必须放在情态小品词 \textit{jo} 前面）
  \item \textit{I morgen når hun \dabold{nok} ikke at sende rapporten.}（V2 倒装后顺序：动词 + 主语 + \textit{nok} + \textit{ikke}）
  \item \textit{Hun blev overrasket, fordi hun \da{jo} ikke kendte reglerne.}（从句中 \textit{jo} 置于主语后、\textit{ikke} 与变位动词前）
  \item \textit{De var \da{da} meget glade for forslaget.}（主句中位于变位动词 \textit{var} 之后）
  \item \textit{Direktøren rejser \dabold{vist} til Berlin på tirsdag.}（位于第一变位动词 \textit{rejser} 之后）
  \item \textit{Han tror, at mødet \dabold{vist} begynder klokken ti.}（\textit{at} 引导从句，\textit{vist} 放在变位动词 \textit{begynder} 前面）
  \item \textit{Du har \da{da} glemt din paraply igen.}（位于助动词 \textit{har} 之后、过去分词 \textit{glemt} 之前）
\end{enumerate}

\textbf{练习 B 答案}
\begin{enumerate}
  \item \dabold{nok}（宽慰、有把握的安慰：“肯定没问题的”）
  \item \da{jo}（提醒双方早已知道的事实：“你毕竟来过三次了呀”）
  \item \dabold{vist}（记忆不确切或模糊估算：“大概五百克朗左右吧”）
  \item \da{da}（表达感叹与共情：“那可真是太遗憾了”）
  \item \dabold{nok} / \dabold{vist}（合理推测他恐怕是忘了，用 \textit{nok} 或表示揣测的 \textit{vist} 均可）
  \item \da{Jo}（注意此空用于反驳否定问句 \textit{Kommer du ikke...?}，须重读用 \textit{Jo}）
\end{enumerate}

\textbf{练习 C 答案}
\begin{enumerate}
  \item \textit{Bare rolig, i morgen kommer han nok ikke for sent.}
  \item \textit{Jeg har jo allerede sendt mailen til dig.}
  \item \textit{Vores chef skal vist til Aarhus i næste uge.}
  \item \textit{Du kan da sagtens stille din kollega det spørgsmål!}
  \item \textit{Hun kom ikke på arbejde, fordi hun vist er syg.}
  \item \textit{Han har nok ret, men vi skal lige tjekke tallene igen.}
\end{enumerate}

\textbf{DU3 Modul 4 高分范文（148 词，不含署名）}\\
\textit{Hej Mikkel}

\textit{Tak for din mail! Jeg kan godt forstå, at du blev lidt usikker efter mødet, men du behøver slet ikke at gå i panik. Danskere bruger nemlig ofte småord, som i starten kan virke lidt forvirrende.}

\textit{Når din chef siger ”det er \dabold{vist} meget fint”, betyder det \da{jo} ikke, at han synes, dit arbejde er dårligt. Ofte bruger danskere bare \dabold{vist}, når de gerne vil kigge lidt nærmere på detaljerne, før de siger endeligt ja. Du er \da{da} lige startet i firmaet, så ingen forventer, at alt er perfekt fra dag ét!}

\textit{Mit bedste råd er, at du spørger ham venligt næste gang: ”Har du eventuelt et par forslag til, hvad jeg kan forbedre i udkastet?” Danske chefer sætter \da{jo} stor pris på åbenhed, og hvis du tager det roligt og spørger nysgerrigt, skal det \dabold{nok} gå rigtig godt mellem jer.}

\textit{Jeg krydser fingre for dig på det nye arbejde!}

\textit{Bedste hilsner\\
\mbox{[Dit navn]}}

\small\textbf{名师点评}：148 词（不含署名）远超 Modul 4 Opgave 1 的 125 词底线。全信将 \textit{vist, jo, da, nok} 自然嵌入主句、否定句与从句，既有情绪共情（\textit{Du er da lige startet...}），又有实操建议与暖心收尾（\textit{skal det nok gå rigtig godt}），语域地道流畅，完全符合私人回复邮件的 B1+/B2- 评分标准。
\end{answerbox}